\phantomsection
\color{unimain}\section*{\centering{\large{ПЕРЕЧЕНЬ СОКРАЩЕНИЙ}}}
\addcontentsline{toc}{section}{Перечень сокращений}
{\color{white}\fontsize{0.1pt}{0.1pt}\selectfont<begABBRS>}
\color{black}
\begin{longtable}{>{\raggedright\arraybackslash}m{2cm}>{\raggedright\arraybackslash}m{0.5cm}>{\raggedright\arraybackslash}m{20cm}}
\endfirsthead\endhead\endfoot\endlastfoot
GOOSE & -- & Generic Object-Oriented Substation Event; \\
IRF & -- & Internal Relay Fault; \\
PPS & -- & Pulse Per Second; \\
PTP & -- & Precision Time Protocol (протокол точного времени); \\
SNTP & -- & Simple Network Time Protocol (простой сетевой протокол времени); \\
АПВ & -- & автоматическое повторное включение; \\
АСУ~ТП & -- & автоматизированная система управления технологическими процессами; \\
АУВ & -- & автоматика управления выключателем; \\
АЦП & -- & аналого-цифровой преобразователь; \\
БИ & -- & блок испытательный; \\
БКУ & -- & блок команд управления; \\
БНН & -- & блокировка при неисправности цепей напряжения; \\
БНТ & -- & блокировка при бросках намагничивающего тока; \\
БСК & -- & батарея статических конденсаторов; \\
БУ & -- & блок управления; \\
ВВ & -- & высоковольтный ввод; \\
ВН & -- & высшее напряжение; \\
ВО & -- & вторичная обмотка; \\
ВЧ & -- & высокочастотный; \\
ВЧПП & -- & высокочастотный приемо-передатчик; \\
ДЗ & -- & дистанционная защита; \\
ДЗТ & -- & дифференциальная защита трансформатора; \\
ДО & -- & дочерние общества; \\
ДОТХ & -- & дифференциальный орган с торможением; \\
ДТЗ & -- & дифференциальная токовая защита; \\
ДТЗт & -- & дифференциальная токовая защита с торможением; \\
ДТО & -- & дифференциальная токовая отсечка; \\
ДФЗ & -- & дифференциально-фазная защита; \\
ЕЭС & -- & единая энергетическая система; \\
ЗАПВ & -- & запрет автоматического повторного включения; \\
ЗВП & -- & защита от внутренних повреждений; \\
ЗИП & -- & запасные части, инструменты и принадлежности; \\
ЗМН & -- & защита минимального напряжения; \\
ЗН & -- & заземляющий нож (разъединитель); \\
ЗНР & -- & защита от неполнофазного режима; \\
ЗНФ & -- & защита от непереключения фаз; \\
ЗПВГ & -- & защита от перегрузки высшими гармониками; \\
ЗПН & -- & защита от повышения напряжения; \\
ИО & -- & измерительный орган / избирательный орган; \\
ИЧМ & -- & интерфейс человек-машина; \\
КЗ & -- & короткое замыкание; \\
КИ & -- & контроль изоляции; \\
КСВ & -- & контроль силового выключателя; \\
КЦТ & -- & контроль цепей тока; \\
ЛВ & -- & логика включения; \\
ЛО & -- & логика отключения; \\
ЛЭП & -- & линия электропередачи; \\
МТЗ & -- & максимальная токовая защита; \\
МЭК & -- & международная электротехническая комиссия; \\
НВ & -- & нейтральный вывод; \\
НВЧЗ & -- & направленная высокочастотная защита; \\
НЗ & -- & нормально закрытый; \\
НО & -- & нормально открытый; \\
НП & -- & нулевая последовательность; \\
ОАПВ & -- & однофазное автоматическое повторное включение; \\
ОВ & -- & оперативный вывод / обходной выключатель; \\
ОЗУ & -- & оперативное запоминающее устройство; \\
ОС & -- & ожидание синхронизма; \\
ОТ & -- & оперативный ток; \\
ПА & -- & противоаварийная автоматика; \\
ПАВ & -- & полуавтоматическое включение; \\
ПЗУ & -- & постоянное запоминающее устройство; \\
ПК & -- & предохранительный клапан; \\
ПО & -- & пусковой орган / программное обеспечение; \\
ПРД & -- & передатчик; \\
ПРМ & -- & приемник; \\
ПУ & -- & пульт управления / поперечное ускорение; \\
РД & -- & реле давления; \\
РЗ & -- & релейная защита; \\
РЗА & -- & релейная защита и автоматика; \\
РКВ & -- & реле команды «включить»; \\
РКО & -- & реле команды «отключить»; \\
РПВ & -- & реле положения «включено»; \\
РС & -- & регистрация событий / реле сопротивления; \\
РТ & -- & реле тока; \\
РУ & -- & распределительное устройство; \\
РФ & -- & реле фиксации; \\
РФК & -- & реле фиксации команд; \\
РФП & -- & реле фиксации включенного положения; \\
РЭ & -- & руководство по эксплуатации; \\
СКРМ & -- & средства компенсации реактивной мощности; \\
СО & -- & система охлаждения; \\
ССПИ & -- & система сбора и передачи информации; \\
ТЗНП & -- & токовая защита нулевой последовательности; \\
ТЗОП & -- & токовая защита обратной последовательности; \\
ТН & -- & трансформатор напряжения; \\
ТР & -- & технический регламент; \\
ТС & -- & технологическая сигнализация; \\
ТТ & -- & трансформатор тока; \\
ТУ & -- & телеуправление / телеускорение / технические условия; \\
УВ & -- & управление выключателем / управляющее воздействие; \\
УРОВ & -- & устройство резервирования при отказе выключателя; \\
ФБ & -- & функциональный блок; \\
ФК & -- & функциональная клавиша; \\
ФПВ & -- & фиксация положения выключателя; \\
ФСУ & -- & функциональная схема устройства; \\
ЦВ & -- & цепи включения; \\
ЦН & -- & цепи напряжения; \\
ЦО & -- & цепи отключения; \\
ЦП & -- & центральный процессор; \\
ЦПС & -- & цифровая подстанция; \\
ЦТ & -- & цепи тока; \\
ЦУ & -- & цепи управления; \\
ШП & -- & шинки питания; \\
ШЭТ & -- & шкаф электротехнический типовой; \\
ЭМВ & -- & электромагнит включения; \\
ЭМО & -- & электромагнит отключения; \\
ЭМУ & -- & электромагнит управления. \\
{\color{white}\fontsize{0.1pt}{0.1pt}\selectfont<endABBRS>}
\end{longtable}
